%============================================================
% CH1 · PATTERN 4 — EXERCISES, HINTS, ANSWERS
%============================================================
\subsection{Pattern 4: Exercises and Going Further}
\label{ssec:ch1-p4-ex}

\begin{patternbox}{Pattern 4 Summary Card}
\begin{tabularx}{\linewidth}{@{}>{\bfseries}l X@{}}
Vertex & a \emph{state}: the smallest snapshot that decides the next moves and the goal test \\
Edge & one move of the statement, generated on the fly by a neighbour function \\
Algorithm & BFS from the start (unit costs); \texttt{dist} as a flat array over integer-encoded states \\
Size & $\#\text{states}\times\#\text{moves}$; unbounded values need a \emph{proven} window \\
Key variations & reverse BFS from the goal, multi-query lookups, path printing, bidirectional BFS, reverse greedy \\
Top traps & no / too small window; marking on pop; forbidden start; negative \texttt{\%}; DFS for ``fewest'' \\
\end{tabularx}
\end{patternbox}

\exlegend
\begin{enumerate}[leftmargin=*, label=\textbf{E4.\arabic*}, itemsep=6pt]
    \setlength{\emergencystretch}{2em}
    \item \exeasy\ \pconcept \cf{1000}\ \exprob{Transformation: from A to B}{Codeforces 727A}{https://codeforces.com/problemset/problem/727/A}. Turn $a$ into $b$ ($1 \le a < b \le 10^9$) using $x \to 2x$ and $x \to 10x + 1$. Print \texttt{NO}, or \texttt{YES}, the length $k$ of a sequence $a = x_1, \dots, x_k = b$ and the sequence itself (any valid one).
    \item \exmedium\ \ptricky \exown{Three buttons.} Two Buttons with a third button that multiplies by 3: moves $x \to 2x$, $x \to 3x$, $x \to x-1$ (the display must stay positive), $1 \le n, m \le 10^4$. Fewest presses from $n$ to $m$. ($5 \to 7$: 3; \ $1 \to 26$: 4.)
    \item \exmedium\ \pclassic \exprob{Water and Jug Problem}{LeetCode 365}{https://leetcode.com/problems/water-and-jug-problem/}\ \lc{Medium}. Jugs of capacities $x$ and $y$ ($\le 10^3$), unlimited water. Moves: fill a jug, empty a jug, pour one into the other until the source is empty or the receiver is full. Can the total amount of water in the two jugs become exactly \texttt{target}?
    \item \exmedium\ \pconcept \cf{1300}\ \exprob{Getting Zero}{Codeforces 1661B}{https://codeforces.com/problemset/problem/1661/B}. One step turns $v$ into $(v+1) \bmod 32768$ or $2v \bmod 32768$. For each of $n \le 32768$ given numbers, print the fewest steps to reach 0.
    \item \exmedium\ \ptricky \exprob{Minimum Operations to Convert Number}{LeetCode 2059}{https://leetcode.com/problems/minimum-operations-to-convert-number/}\ \lc{Medium}. From $x = \texttt{start}$ ($0 \le \texttt{start} \le 1000$), one operation replaces $x$ by $x + a_i$, $x - a_i$ or $x \oplus a_i$ for some $a_i$ of the array, but only while $0 \le x \le 1000$; results outside the range are allowed but final. Fewest operations to obtain \texttt{goal}, or $-1$.
    \item \exmedium\ \precent \exprob{Minimum Number of Operations to Make X and Y Equal}{LeetCode 2998}{https://leetcode.com/problems/minimum-number-of-operations-to-make-x-and-y-equal/}\ \lc{Medium}. One operation on $x$: divide by 11 (if divisible), divide by 5 (if divisible), $-1$ or $+1$. Fewest operations to turn $x$ into $y$ ($1 \le x, y \le 10^4$).
\end{enumerate}

\subsubsection*{Hints}
\begin{enumerate}[leftmargin=*, label=\textbf{E4.\arabic*}]
    \item Both moves only increase the number. Look at the \emph{last} digit of $b$: which move could have produced it?
    \item Repeat the Lemma of Problem~1.4.1 for any multiplication: what is the window now? Is the reverse greedy still safe?
    \item State $=$ the pair of amounts $(a, b)$. There are only $(x{+}1)(y{+}1)$ of them. Bonus: which totals are reachable at all (think of $\gcd(x,y)$)?
    \item $n$ separate BFS runs are $32768^2$ steps. All queries share the same goal --- search backwards from 0 once.
    \item The vertices are $0..1000$. A generated value outside the range is not a vertex, but it may still be the goal.
    \item $+1$ can push $x$ \emph{above} both $x$ and $y$ (to reach a multiple of 11). How far up can an optimal path go, if the answer is never more than $|x - y|$?
\end{enumerate}

\subsubsection*{Answers}

\begin{solutionbox}[E4.1 Transformation: from A to B]
\textbf{Model:} vertices are the numbers $a..b$, edges $x \to 2x$ and $x \to 10x+1$; both increase, so nothing above $b$ matters. \textbf{Backwards} every number has at most one predecessor: an even number can only come from $\times2$ (since $10x + 1$ is odd), a number ending in 1 only from appending 1, any other odd number has none. So the reverse walk is forced --- a path in a graph where every vertex has out-degree at most one (Pattern~5 flavour).
\begin{lstlisting}[style=cpp]
#include <bits/stdc++.h>
using namespace std;

int main() {
    long long a, b;
    cin >> a >> b;
    vector<long long> path = {b};
    while (b > a) {
        if (b % 2 == 0) b /= 2;                 // undo x -> 2x
        else if (b % 10 == 1) b /= 10;          // undo x -> 10x + 1
        else break;                             // no predecessor
        path.push_back(b);
    }
    if (b != a) { cout << "NO\n"; return 0; }
    reverse(path.begin(), path.end());
    cout << "YES\n" << path.size() << "\n";
    for (long long v : path) cout << v << " ";
    cout << "\n";
}
\end{lstlisting}
\begin{center}
\begin{tikzpicture}[node distance=1.35cm]
    \node[gvdst, minimum size=10mm, font=\scriptsize\bfseries] (a) {162};
    \node[gv, minimum size=10mm, font=\scriptsize\bfseries, right=of a] (b) {81};
    \node[gv, minimum size=10mm, font=\scriptsize\bfseries, right=of b] (c) {8};
    \node[gv, minimum size=10mm, font=\scriptsize\bfseries, right=of c] (d) {4};
    \node[gvsrc, minimum size=10mm, font=\scriptsize\bfseries, right=of d] (e) {2};
    \draw[back] (a) -- node[lbl, above]{$\div2$} (b);
    \draw[back] (b) -- node[lbl, above]{drop 1} (c);
    \draw[back] (c) -- node[lbl, above]{$\div2$} (d);
    \draw[back] (d) -- node[lbl, above]{$\div2$} (e);
\end{tikzpicture}
\end{center}
Samples: $(2, 162)$ $\to$ \texttt{YES}, 5, \texttt{2 4 8 81 162}; \ $(4, 42)$ $\to$ \texttt{NO} ($42 \to 21 \to 2 < 4$); \ $(100, 40021)$ $\to$ \texttt{YES}, 5, \texttt{100 200 2001 4002 40021}. Time $\bigO{\log b}$. A forward BFS over $a..b$ would also be correct but has up to $10^9$ vertices.
\end{solutionbox}

\begin{solutionbox}[E4.2 Three buttons]
\textbf{Window.} The Lemma's proof never used the factor 2: a multiplication made at $x \ge m$ adds at least $x \ge 1$ per press that must be removed by blue presses, so it is never optimal. Multiplications happen only at $x < m$, giving at most $3(m-1)$: $\text{HI} = 3\max(n, m)$.
\begin{lstlisting}[style=cpp]
#include <bits/stdc++.h>
using namespace std;

int main() {
    int n, m;
    cin >> n >> m;
    const int HI = 3 * max(n, m);
    vector<int> dist(HI + 1, -1);
    queue<int> q;
    dist[n] = 0;
    q.push(n);
    while (!q.empty()) {
        int x = q.front(); q.pop();
        if (x == m) break;
        int next[3] = {2 * x, 3 * x, x - 1};
        for (int y : next)
            if (y >= 1 && y <= HI && dist[y] == -1) { dist[y] = dist[x] + 1; q.push(y); }
    }
    cout << dist[m] << "\n";
}
\end{lstlisting}
$5 \to 7$: 3 ($5\to4\to8\to7$); \ $4 \to 11$: 2 ($4\to12\to11$); \ $1 \to 26$: 4 ($1\to3\to9\to27\to26$). The tempting reverse greedy ``divide by 3 if possible, else by 2 if even, else add 1'' gives 7 on $1 \to 26$ ($26\to13\to14\to7\to8\to4\to2\to1$): with two multiplications the backward choice is no longer forced, and BFS is the safe method. (Checked: BFS with $\text{HI} = 3\max(n,m)$ equals BFS on $[1, 10^5]$ for all $n, m \le 150$.)
\end{solutionbox}

\begin{solutionbox}[E4.3 Water and Jug Problem]
\textbf{Model:} state $(a, b)$, $0 \le a \le x$, $0 \le b \le y$, id $a(y{+}1) + b$; six moves; goal test ``$a + b = \texttt{target}$'' on every popped state.
\begin{lstlisting}[style=cpp]
#include <bits/stdc++.h>
using namespace std;

class Solution {
public:
    bool canMeasureWater(int x, int y, int target) {
        if (target > x + y) return false;
        vector<char> seen((x + 1) * (y + 1), 0);
        queue<pair<int, int>> q;
        q.push(make_pair(0, 0));
        seen[0] = 1;
        while (!q.empty()) {
            int a = q.front().first, b = q.front().second;
            q.pop();
            if (a + b == target) return true;
            int ab = min(a, y - b), ba = min(b, x - a);       // amounts poured
            int nxt[6][2] = {{x, b}, {a, y}, {0, b}, {a, 0},  // fill / empty
                             {a - ab, b + ab}, {a + ba, b - ba}};  // pour a->b, b->a
            for (auto& s : nxt) {
                int id = s[0] * (y + 1) + s[1];
                if (!seen[id]) { seen[id] = 1; q.push(make_pair(s[0], s[1])); }
            }
        }
        return false;
    }
};
\end{lstlisting}
At most $1001^2 \approx 10^6$ states $\times$ 6 moves: fast. Samples: $(3,5,4) \to$ \texttt{true}, $(2,6,5) \to$ \texttt{false}, $(1,2,3) \to$ \texttt{true}.

\begin{center}
\begin{tikzpicture}[node distance=0.9cm]
    \foreach \s/\lab [count=\i from 0] in {gvsrc/{0,0}, gv/{0,5}, gv/{3,2}, gv/{0,2}, gv/{2,0}, gv/{2,5}, gv/{3,4}, gvdst/{0,4}}
        \node[\s, minimum size=10mm, font=\scriptsize\bfseries] (j\i) at (\i*1.65,0) {\lab};
    \foreach \i [evaluate=\i as \j using int(\i+1)] in {0,...,6} \draw[gpath, -{Stealth[length=2.2mm]}] (j\i) -- (j\j);
\end{tikzpicture}
\end{center}
For $x=3$, $y=5$ only 16 of the 24 states are reachable, and the first state with total 4 is at distance 7, e.g.\ the path above (also $(3,1)$ at distance 7).

\textbf{Mathematical shortcut.} Every reachable amount is a combination $px + qy$, hence a multiple of $g = \gcd(x, y)$, and conversely every multiple of $g$ up to $x + y$ is reachable (B\'ezout). So the answer is $\texttt{target} \le x + y$ and $g \mid \texttt{target}$ --- checked against the BFS for all $x, y \le 40$, $\texttt{target} \le 90$.
\end{solutionbox}

\begin{solutionbox}[E4.4 Getting Zero]
\textbf{Model:} $32768$ residues, edges $v \to v+1$ and $v \to 2v$ (mod $32768$). Every query has the same goal 0, so run \textbf{one BFS from 0 on the reversed edges}: the predecessors of $v$ are $v - 1$ and, if $v$ is even, the two solutions of $2u \equiv v$: $u = v/2$ and $u = v/2 + 16384$.

\begin{center}
\begin{tikzpicture}[node distance=1.4cm]
    \node[gvsrc, minimum size=11mm, font=\scriptsize\bfseries] (z) {0};
    \node[gvseen, minimum size=11mm, font=\scriptsize\bfseries, right=3.2cm of z, yshift=1.6cm] (a) {32767};
    \node[gvseen, minimum size=11mm, font=\scriptsize\bfseries, right=3.2cm of z] (b) {16384};
    \node[gvdead, minimum size=11mm, font=\scriptsize\bfseries, right=3.2cm of z, yshift=-1.6cm] (c) {0};
    \draw[gde] (z) -- node[gw, above, sloped]{undo $+1$} (a);
    \draw[gde] (z) -- node[gw, above]{undo $\times2$} (b);
    \draw[gfaded, dashed, -{Stealth[length=2mm]}] (z) -- node[gw, below, sloped]{$0/2$} (c);
    \node[note, anchor=west, text width=5cm, align=left] at (6.0,0) {reverse BFS from 0: layer 1 is $\{32767, 16384\}$; afterwards every query is \texttt{dist[a\_i]}.};
\end{tikzpicture}
\end{center}
\begin{lstlisting}[style=cpp]
#include <bits/stdc++.h>
using namespace std;

int main() {
    const int M = 32768;
    vector<int> dist(M, -1);
    queue<int> q;
    dist[0] = 0;
    q.push(0);
    while (!q.empty()) {
        int v = q.front(); q.pop();
        vector<int> pre = {(v - 1 + M) % M};          // u + 1 = v
        if (v % 2 == 0) { pre.push_back(v / 2); pre.push_back(v / 2 + M / 2); }  // 2u = v
        for (int u : pre)
            if (dist[u] == -1) { dist[u] = dist[v] + 1; q.push(u); }
    }
    int n;
    cin >> n;
    for (int i = 0; i < n; i++) {
        int a;
        cin >> a;
        cout << dist[a] << (i + 1 < n ? ' ' : '\n');
    }
}
\end{lstlisting}
Sample: $19, 32764, 10240, 49 \to 14\ 4\ 4\ 15$. Time $\bigO{M + n}$. \textbf{Alternative:} since $15$ doublings always reach 0, the answer is at most 15, so it is $\min$ over $\text{add} = 0..15$ of $\text{add} + $ (doublings needed for $v + \text{add}$, which is $15 - $ the number of trailing zero bits); both methods agree on all $32768$ inputs, the maximum being 15.
\end{solutionbox}

\begin{solutionbox}[E4.5 Minimum Operations to Convert Number]
\textbf{Model:} vertices $0..1000$; from $x$ there are $3|\texttt{nums}|$ edges. Values outside $[0, 1000]$ are leaves: compare them with \texttt{goal} \emph{before} the range check (the goal may be $-4$ or $10^9$).
\begin{lstlisting}[style=cpp]
#include <bits/stdc++.h>
using namespace std;

class Solution {
public:
    int minimumOperations(vector<int>& nums, int start, int goal) {
        vector<int> dist(1001, -1);
        queue<int> q;
        dist[start] = 0;
        q.push(start);
        while (!q.empty()) {
            int x = q.front(); q.pop();
            for (int a : nums) {
                long long cand[3] = {(long long)x + a, (long long)x - a, (long long)(x ^ a)};
                for (long long y : cand) {
                    if (y == goal) return dist[x] + 1;         // goal may be out of range
                    if (y < 0 || y > 1000 || dist[y] != -1) continue;
                    dist[y] = dist[x] + 1;
                    q.push((int)y);
                }
            }
        }
        return -1;
    }
};
\end{lstlisting}
Returning when the goal is \emph{generated} is safe here: BFS generates vertices in non-decreasing distance, so the first generation of the goal uses the fewest operations (the goal is never pushed, so the ``popped'' test would miss out-of-range goals). Samples: $([2,4,12], 2, 12) \to 2$, \ $([3,5,7], 0, -4) \to 2$, \ $([2,8,16], 0, 1) \to -1$. Time $\bigO{1001 \cdot 3|\texttt{nums}|} \approx 3\cdot10^6$.
\end{solutionbox}

\begin{solutionbox}[E4.6 Minimum Number of Operations to Make X and Y Equal]
\textbf{Model:} vertices are numbers, edges $v \to v-1$, $v \to v+1$, $v \to v/11$ and $v \to v/5$ (when divisible); BFS from $x$ to $y$. \textbf{Window.} Stepping by $\pm1$ alone costs $|x - y|$, so the answer is at most $|x - y| < M = \max(x, y)$. Only $+1$ increases the value, so an optimal path never climbs above $x + |x - y| \le 2M$; and 0 is the smallest value one can reach. BFS on $[0, 2M]$ is exact. The window \emph{must} go above $\max(x,y)$: in sample 2, $54 \to 55$ is the first move.

\begin{center}
\begin{tikzpicture}[node distance=1.5cm]
    \node[gvsrc] (a) {54};
    \node[gv, right=of a] (b) {55};
    \node[gv, right=of b] (c) {5};
    \node[gv, right=of c] (d) {1};
    \node[gvdst, right=of d] (e) {2};
    \draw[gpath, -{Stealth[length=2.4mm]}] (a) -- node[gw, above=1pt]{$+1$} (b);
    \draw[gpath, -{Stealth[length=2.4mm]}] (b) -- node[gw, above=1pt]{$\div11$} (c);
    \draw[gpath, -{Stealth[length=2.4mm]}] (c) -- node[gw, above=1pt]{$\div5$} (d);
    \draw[gpath, -{Stealth[length=2.4mm]}] (d) -- node[gw, above=1pt]{$+1$} (e);
\end{tikzpicture}
\end{center}
\begin{lstlisting}[style=cpp]
#include <bits/stdc++.h>
using namespace std;

class Solution {
public:
    int minimumOperationsToMakeEqual(int x, int y) {
        const int HI = 2 * max(x, y);                   // proven window [0, HI]
        vector<int> dist(HI + 1, -1);
        queue<int> q;
        dist[x] = 0;
        q.push(x);
        while (!q.empty()) {
            int v = q.front(); q.pop();
            if (v == y) return dist[v];
            vector<int> nxt = {v - 1, v + 1};
            if (v % 11 == 0) nxt.push_back(v / 11);
            if (v % 5 == 0) nxt.push_back(v / 5);
            for (int w : nxt)
                if (w >= 0 && w <= HI && dist[w] == -1) { dist[w] = dist[v] + 1; q.push(w); }
        }
        return -1;                                      // unreachable: +1/-1 connect all
    }
};
\end{lstlisting}
Samples: $(26, 1) \to 3$, \ $(54, 2) \to 4$, \ $(25, 30) \to 5$. With $x \le y$ no division helps (it only lowers the value), and the answer is $y - x$. Checked: the window $[0, 2M]$ gives the same answers as $[0, 10^5]$ for all $x, y \le 300$ and 300 random pairs up to $10^4$, whereas the window $[0, M]$ fails on 30382 pairs with $x \le 2000$, $y \le 50$. Time $\bigO{M}$.
\end{solutionbox}
